\documentclass{article}
\usepackage{tabularx}
\usepackage[table]{xcolor}
\usepackage{graphicx}
\usepackage{datetime}

\begin{document}

    % reduce white space at beginning
    \title{\vspace{-6cm}}
    \date{}
    \maketitle
    % do not print a page number
    \thispagestyle{empty}


    \begin{center}
        \huge \VAR[title|latex]
    \end{center}

    % Rechnungsnummer und Datum
    \vspace{1cm}
    \textit{TEST \today - \currenttime } % for testing, to see if a new bon was created

    \textbf{Rechnung Nr. \VAR["{:010}".format(order.id)]} \hspace{\fill}  \VAR[order.finished_at.date()]

    % Addresse und Steuer ID
    \begin{flushright}
        \VAR[issuer|latex] \\
        % \VAR[address.strip().replace("\n", "\\\\")|latex] \\
        \VAR[address.strip()|latex] \\
        USt-IdNr. \VAR[ust_id|latex] \\
    \end{flushright}

    % Line Items in the Rechnung
    \vspace{1cm}
    \newcolumntype{R}{>{\raggedleft\arraybackslash}X}
    \rowcolors{2}{}{gray!40}
    \begin{tabularx}{\textwidth}{ p{1cm} p{4cm} R R X }
        ~                    & \textbf{Produkt}        & \textbf{Einzelpreis}   & \textbf{Gesamt}              &                       \\
        \hline
        \BLOCK[ for item in order.line_items ]
        \VAR[item.quantity]x & \VAR[item.product.name] & \VAR[item.price|money] & \VAR[item.total_price|money] & [\VAR[item.tax_name]] \\
        \BLOCK[endfor]
    \end{tabularx}


    % Gesamtsumme
    \vspace{1cm}
    \hspace{\fill} \textbf{Summe EUR} \textbf{\VAR[order.value_sum|money]€}

    %Steuerstatz
    \vspace{1cm}
    \rowcolors{0}{}{}
    \begin{tabularx}{\textwidth}{ RRRR }
        \textbf{MwSt.\%}                                         & \textbf{Brutto}                & \textbf{Netto}                   & \textbf{MwSt.}                 \\
        \hline
        \BLOCK[ for tax_rate in tax_rates ]
        \VAR[tax_rate.tax_name]=\VAR[ tax_rate.tax_rate|percent] & \VAR[tax_rate.value_sum|money] & \VAR[tax_rate.value_notax|money] & \VAR[tax_rate.value_tax|money] \\
        \BLOCK[endfor]
        \hline
        ~                                                        & \VAR[order.value_sum|money]    & \VAR[order.value_notax|money]    & \VAR[order.value_tax|money]    \\
    \end{tabularx}

    % sonstige Infos
    \vspace{1cm}
    \begin{tabularx}{0.8\textwidth}{ RX }
        Zahlmethode:        & \VAR[order.txmethod]    \\
        Vorgangsbeginn:     & \VAR[order.created_at.astimezone()]  \\
        Vorgangsende:       & \VAR[order.finished_at.astimezone()] \\
        Kasse:              & Festzelt 1              \\
        Kassierer:          & 1337                    \\
        TSE Seriennummer:   & asdfqwer                \\
        TSE Signaturzähler: & signatur                \\
        TSE Prüfwert:       & prüfwert                \\
    \end{tabularx}

    % evtl direkter online Link um die Rechnung abzurufen

    % ein kurzer "witz"
    \vspace{\fill}
    \begin{center}
        \VAR[funny_text|latex]

        \includegraphics[width=3cm]{logo}
    \end{center}

\end{document}
